\chapter{Noetherian Rings and Hilbert's Theorems}

\begin{center}\itshape
Can every system of polynomial equations, even an infinite one, be replaced by a finite one? When does a system have no solution at all? Hilbert's basis theorem and his
Nullstellensatz answer both questions, turn ideals into geometry, and, through Gröbner bases, into algorithms.
\end{center}

In the 1880s the theory of invariants was the most fashionable subject in algebra, and Paul Gordan of Erlangen was its king. A binary form, such as $ax^2+bxy+cy^2$, can be
transformed by a linear change of the variables, and an \emph{invariant} is a polynomial in the coefficients that does not change, like the discriminant $b^2-4ac$. The central
question was whether all invariants of a given kind are polynomials in finitely many of them. Gordan had answered it in 1868 for binary forms, by computations so long and
ingenious that nobody could extend them to more variables.

In 1888 a young Privatdozent from Königsberg, David Hilbert, settled the question for all forms in a few pages, and he did it without computing a single invariant. He proved that
in any polynomial ring, every ideal is generated by finitely many elements; finiteness of invariants followed by an argument by contradiction, which showed that a finite set of
generators must exist without saying how to find it. Gordan's reaction is famous, and comes to us through Hilbert's biographer Constance Reid: ``This is not mathematics. This is
theology.'' Felix Klein saw at once the importance of the method, Hilbert's paper appeared in the \emph{Mathematische Annalen} in 1890, and Gordan himself later conceded that
theology also has its merits. Three years later Hilbert proved his second great theorem on polynomial ideals, the Nullstellensatz, or ``theorem of the zeros''.

This chapter proves both theorems, and draws the dictionary between algebra and geometry that they make possible: ideals and the sets where polynomials vanish, prime ideals and
irreducible sets, maximal ideals and points. Volume~6, Chapter~3 promised the topology of algebraic geometry; here it is. The chapter ends with the algorithmic answer to Gordan's
objection: Gröbner bases, which Bruno Buchberger invented in his 1965 thesis and which turn Hilbert's existence proofs into computations.

All rings are commutative with $1$.

\section{Noetherian rings and modules}

\begin{proposition}\label{prop:noetherian}
For a ring $R$, the following are equivalent:
\begin{enumerate}[label=(\alph*)]
\item every ideal of $R$ is finitely generated;
\item every increasing chain of ideals $I_1\subseteq I_2\subseteq\cdots$ is eventually constant (the \emph{ascending chain condition});
\item every non-empty set of ideals has a maximal element for inclusion.
\end{enumerate}
A ring with these properties is called \emph{Noetherian}, after Emmy Noether, who made the ascending chain condition the foundation of the theory in 1921. The same equivalences, with the same
proof, hold for the submodules of a module, and define \emph{Noetherian modules}.
\end{proposition}

\begin{proof}
(a)$\Rightarrow$(b). The union $I$ of a chain of ideals is an ideal: two of its elements lie in a common $I_n$. If $I=(a_1,\dots,a_k)$, each $a_i$ lies in some $I_{n_i}$, and for $n\ge\max n_i$, $I\subseteq I_n\subseteq I$.
(b)$\Rightarrow$(c). If a non-empty set $\Sigma$ of ideals had no maximal element, we could choose $I_1\in\Sigma$, then $I_2\in\Sigma$ with $I_1\subsetneq I_2$, and so on, an increasing chain that never stops.
(c)$\Rightarrow$(a). Let $I$ be an ideal and $\Sigma$ the set of finitely generated ideals contained in $I$; it contains $(0)$. Let $J$ be maximal in $\Sigma$. For $a\in I$, $J+(a)$ is in $\Sigma$ and contains $J$, so equals $J$: $a\in J$.
Hence $I=J$ is finitely generated.
\end{proof}

Fields and principal ideal domains are Noetherian, $\Z$ and $K[X]$ in particular. The ring $K[X_1,X_2,\dots]$ of polynomials in infinitely many variables is not: the chain
$(X_1)\subsetneq(X_1,X_2)\subsetneq\cdots$ never stops. By (c), in a Noetherian ring every proper ideal $I$ is contained in a maximal ideal: take a maximal element of the set of proper ideals containing $I$.

\begin{proposition}\label{prop:noethmodules}
Let $0\to M'\to M\to M''\to0$ be a short exact sequence of modules. Then $M$ is Noetherian if and only if $M'$ and $M''$ are. Consequently, over a Noetherian ring every finitely generated module is
Noetherian, and so is each of its submodules.
\end{proposition}

\begin{proof}
Identify $M'$ with a submodule of $M$ and $M''$ with $M/M'$. Submodules of $M'$ are submodules of $M$, and submodules of $M/M'$ are the $N/M'$ for submodules $N\supseteq M'$ (Volume~2, Chapter~7), so chains in
$M'$ or $M''$ come from chains in $M$: if $M$ is Noetherian, so are $M'$ and $M''$. Conversely let $N_1\subseteq N_2\subseteq\cdots$ be a chain in $M$. The chains $N_k\cap M'$ and $(N_k+M')/M'$ are eventually constant, say from $k_0$ on. For
$k\ge k_0$ and $x\in N_{k+1}$, there is $y\in N_k$ with $x-y\in M'$, by the second chain; then $x-y\in N_{k+1}\cap M'=N_k\cap M'$, so $x\in N_k$. Hence $N_{k+1}=N_k$.

By induction on $n$, using $0\to R\to R^n\to R^{n-1}\to0$, the module $R^n$ is Noetherian when $R$ is. A finitely generated module is a quotient of some $R^n$.
\end{proof}

\section{Hilbert's basis theorem}

\begin{theorem}[Hilbert's basis theorem, 1890]\label{thm:hilbertbasis}
If $R$ is Noetherian, so is $R[X]$. Hence $K[X_1,\dots,X_n]$ and $\Z[X_1,\dots,X_n]$ are Noetherian, for every field $K$, and so is every finitely generated algebra over a Noetherian ring.
\end{theorem}

\begin{proof}
Let $I$ be an ideal of $R[X]$. For $d\ge0$ let $L_d$ be the set of leading coefficients of the elements of $I$ of degree $d$, together with $0$. It is an ideal of $R$: if $f,g\in I$ have degree $d$ and leading coefficients
$a,b$, then $f+g$ and $rf$ have leading coefficients $a+b$ and $ra$, or are of lower degree when these vanish. Moreover $L_d\subseteq L_{d+1}$, because $Xf\in I$ has degree $d+1$ and the same leading coefficient. By the ascending
chain condition, $L_d=L_N$ for all $d\ge N$. Each of $L_0,\dots,L_N$ is finitely generated; for each $d\le N$ choose polynomials $f_{d,1},\dots,f_{d,k_d}\in I$ of degree $d$ whose leading coefficients generate $L_d$. Let $J$ be the ideal they
generate; we show by induction on the degree that every $f\in I$ lies in $J$.

If $f=0$ there is nothing to prove. Let $f\in I$ have degree $d$ and leading coefficient $a\in L_d$, and let $e=\min(d,N)$, so that $a\in L_e$. Write $a=\sum_jr_j\operatorname{lc}(f_{e,j})$. Then
\[g=f-\sum_jr_jX^{d-e}f_{e,j}\]
lies in $I$ and has degree less than $d$, since the coefficients of $X^d$ cancel. By induction $g\in J$, so $f\in J$. The base case is covered by the same argument: a constant $f\in I$ has degree $0$, and $g=0$.

By induction on $n$, $R[X_1,\dots,X_n]=R[X_1,\dots,X_{n-1}][X_n]$ is Noetherian. A finitely generated $R$-algebra is a quotient of some $R[X_1,\dots,X_n]$, and the ideals of a quotient correspond to ideals containing the kernel.
\end{proof}

The proof is Hilbert's argument in modern dress, and it shows what upset Gordan: the generators are obtained from the ascending chain condition, which says that a chain stops without saying when.

\begin{corollary}\label{cor:finitesystem}
Every system of polynomial equations $\{f=0:f\in S\}$ in $n$ unknowns, with $S\subseteq K[X_1,\dots,X_n]$ arbitrary, has the same solutions as a finite subsystem.
\end{corollary}

\begin{proof}
The ideal $(S)$ is generated by finitely many elements, each a combination of finitely many elements of $S$; collect these into a finite $S_0\subseteq S$. Then $(S_0)=(S)$, and a point where all of $S_0$ vanishes
is a point where every element of $(S_0)$ vanishes, hence all of $S$.
\end{proof}

\section{Varieties and the Zariski topology}

From now on $K$ is a field and $A=K[X_1,\dots,X_n]$. For $S\subseteq A$, the \emph{zero set} of $S$ is
\[V(S)=\{x\in K^n:f(x)=0\text{ for all }f\in S\},\]
and for $Y\subseteq K^n$, the \emph{ideal} of $Y$ is $I(Y)=\{f\in A:f(y)=0\text{ for all }y\in Y\}$. Zero sets are called \emph{algebraic sets}, or \emph{affine varieties}. Clearly $V(S)=V((S))$, both $V$ and $I$ reverse inclusions,
$Y\subseteq V(I(Y))$ and $J\subseteq I(V(J))$.

\begin{proposition}[the Zariski topology]\label{prop:zariski}
$V(A)=\varnothing$, $V(0)=K^n$, $V(I)\cup V(J)=V(IJ)$ and $\bigcap_\lambda V(I_\lambda)=V\bigl(\sum_\lambda I_\lambda\bigr)$. Hence the algebraic sets are the closed sets of a topology on $K^n$, the \emph{Zariski topology}. For a closed set $Y$,
$V(I(Y))=Y$.
\end{proposition}

\begin{proof}
Only the union needs a word. $IJ\subseteq I$ and $IJ\subseteq J$ give $V(I)\cup V(J)\subseteq V(IJ)$. If $x\notin V(I)\cup V(J)$, pick $f\in I$, $g\in J$ with $f(x)\ne0\ne g(x)$; then $fg\in IJ$ and $(fg)(x)\ne0$, so $x\notin V(IJ)$. For the
last statement, if $Y=V(J)$ then $J\subseteq I(Y)$, so $V(I(Y))\subseteq V(J)=Y$; the other inclusion always holds.
\end{proof}

On the line $K$, a non-zero polynomial has finitely many roots, so the closed sets are $K$ and the finite sets: the Zariski topology on $K$ is the cofinite topology of Volume~6, Chapter~3. It is not
Hausdorff when $K$ is infinite, and on $K^2$ it is not the product topology (Exercise~\ref{ex:notproduct}). It is coarse, but it has a property the usual topologies lack.

\begin{definition}\label{def:irreducible}
A non-empty topological space $Y$ is \emph{irreducible} if it is not the union of two closed proper subsets; equivalently, any two non-empty open subsets of $Y$ meet. A space is \emph{Noetherian} if every decreasing
sequence of closed subsets is eventually constant.
\end{definition}

\begin{proposition}\label{prop:components}
\leavevmode
\begin{enumerate}[label=(\alph*)]
\item $K^n$ with its Zariski topology, and each of its subsets, is a Noetherian space.
\item A closed set $Y\subseteq K^n$ is irreducible if and only if $I(Y)$ is a prime ideal.
\item Every non-empty closed subset of a Noetherian space is a finite union $Y=Y_1\cup\dots\cup Y_r$ of irreducible closed subsets, and this decomposition is unique if one requires $Y_i\not\subseteq Y_j$ for $i\ne j$. The
$Y_i$ are the \emph{irreducible components} of $Y$.
\end{enumerate}
\end{proposition}

\begin{proof}
(a) If $Y_1\supseteq Y_2\supseteq\cdots$ are closed in $K^n$, then $I(Y_1)\subseteq I(Y_2)\subseteq\cdots$ stops by Theorem~\ref{thm:hilbertbasis}, and $Y_k=V(I(Y_k))$ stops too. The closed sets of a subspace $Z$ are the
$Y\cap Z$ with $Y$ closed; a decreasing chain $C_k$ of them is $C_k=\overline{C_k}\cap Z$, and the closures $\overline{C_k}$ in $K^n$ decrease, hence stop.

(b) Suppose $Y$ irreducible and $fg\in I(Y)$. Then $Y=(Y\cap V(f))\cup(Y\cap V(g))$, so $Y=Y\cap V(f)$ or $Y=Y\cap V(g)$: $f\in I(Y)$ or $g\in I(Y)$. And $I(Y)\ne A$ since $Y\ne\varnothing$. Conversely, let $I(Y)$ be prime and $Y=Y_1\cup Y_2$ with
$Y_1,Y_2$ closed and proper. Then $I(Y_i)\supsetneq I(Y)$ (as $Y_i=V(I(Y_i))\ne Y$), so there are $f_i\in I(Y_i)\setminus I(Y)$. But $f_1f_2$ vanishes on $Y_1\cup Y_2=Y$, contradicting primality.

(c) Let $\Sigma$ be the set of non-empty closed sets that are not finite unions of irreducible closed sets. If $\Sigma\neq\varnothing$, it has a minimal element $Y$, because the space is Noetherian (a set with no minimal element
would give a strictly decreasing sequence, as in Proposition~\ref{prop:noetherian}). $Y$ is not irreducible, so $Y=Y'\cup Y''$ with $Y',Y''$ closed and proper; by minimality both are finite unions of irreducibles, and so is $Y$, a
contradiction. Remove from a decomposition the $Y_i$ contained in others. For uniqueness, if also $Y=Z_1\cup\dots\cup Z_s$, then $Y_1=\bigcup_j(Y_1\cap Z_j)$, so $Y_1\subseteq Z_j$ for some $j$ by irreducibility, and likewise
$Z_j\subseteq Y_k$; then $Y_1\subseteq Y_k$ forces $k=1$ and $Y_1=Z_j$. Continue with the other $Y_i$.
\end{proof}

\begin{figure}[ht]
\centering
\fig{ch03-varieties}
\caption{Left: the nodal cubic $y^2=x^3+x^2$ in $\R^2$. Its ideal is prime, so it is irreducible, although it crosses itself at the origin. Right: the zero set of $xz$ and $yz$ in $\R^3$, the union of the plane $z=0$ and
the $z$-axis, two irreducible components of different dimensions.}
\label{fig:varieties}
\end{figure}

For instance $V(xy)\subseteq K^2$ is the union of the two axes $V(x)$ and $V(y)$, its irreducible components, and $V(xz,yz)\subseteq K^3$ is the union of the plane $V(z)$ and the line $V(x,y)$ (Figure~\ref{fig:varieties}).

\section{The Nullstellensatz}

Over $\R$, the polynomial $X^2+1$ generates a proper ideal and has no zero, so a proper ideal can have an empty zero set. Hilbert's theorem says that this cannot happen over an algebraically closed field, and the key
is an algebraic lemma. We first need a finiteness result of Emil Artin and John Tate.

\begin{lemma}[Artin and Tate]\label{lem:artintate}
Let $A\subseteq B\subseteq C$ be rings, with $A$ Noetherian, $C$ finitely generated as an $A$-algebra, and $C$ finitely generated as a $B$-module. Then $B$ is finitely generated as an $A$-algebra.
\end{lemma}

\begin{proof}
Let $C=A[c_1,\dots,c_m]$ and $C=By_1+\dots+By_k$. Write $c_i=\sum_jb_{ij}y_j$ and $y_iy_j=\sum_lb_{ijl}y_l$ with coefficients in $B$, and let $B_0=A[b_{ij},b_{ijl}]\subseteq B$. Every element of $C$ is a polynomial in the $c_i$ with
coefficients in $A$; substituting the first relations and then reducing products of the $y$'s with the second, we see that $C=B_0+B_0y_1+\dots+B_0y_k$. So $C$ is a finitely generated $B_0$-module. Now $B_0$ is Noetherian by
Theorem~\ref{thm:hilbertbasis}, so its submodule $B$ is finitely generated by Proposition~\ref{prop:noethmodules}: $B=B_0z_1+\dots+B_0z_p$. Then $B=A[b_{ij},b_{ijl},z_1,\dots,z_p]$.
\end{proof}

\begin{theorem}[Zariski's lemma]\label{thm:zariskilemma}
Let $K$ be a field and $L$ a field that is finitely generated as a $K$-algebra. Then $L$ is a finite extension of $K$.
\end{theorem}

\begin{proof}
Let $L=K[x_1,\dots,x_n]$. If every $x_i$ is algebraic over $K$, then $L$ is obtained by adjoining finitely many algebraic elements, and is a finite extension (Volume~2, Chapter~8). Otherwise, reorder the $x_i$ so that
$x_1,\dots,x_r$ ($r\ge1$) are algebraically independent over $K$ and maximal with this property among the $x_i$. Each $x_j$, $j>r$, is then algebraic over $F=K(x_1,\dots,x_r)\subseteq L$, so $L=F[x_{r+1},\dots,x_n]$ is a finite extension
of $F$, in particular a finitely generated $F$-module. By Lemma~\ref{lem:artintate} with $A=K$, $B=F$, $C=L$, the field $F$ is a finitely generated $K$-algebra: $F=K[p_1/q_1,\dots,p_s/q_s]$ with $p_i,q_i$ in the polynomial ring
$P=K[x_1,\dots,x_r]$. Let $q=q_1\cdots q_s$. Every element of $F$ is then of the form $p/q^k$ with $p\in P$.

Now $P$ is a unique factorisation domain (Volume~2, Theorem~6.17) with infinitely many pairwise non-associate irreducible elements: if $\pi_1,\dots,\pi_t$ were all of them, an irreducible factor of $1+x_1\pi_1\cdots\pi_t$,
which is not a constant, would divide none of them. Choose an irreducible $\pi$ that does not divide $q$. Then $1/\pi=p/q^k$ gives $q^k=\pi p$, so $\pi$ divides $q^k$, hence $q$: a contradiction.
\end{proof}

\begin{theorem}[weak Nullstellensatz]\label{thm:weaknull}
Let $K$ be algebraically closed and $I$ a proper ideal of $K[X_1,\dots,X_n]$. Then $V(I)\ne\varnothing$. Moreover the maximal ideals of $K[X_1,\dots,X_n]$ are exactly the ideals $(X_1-a_1,\dots,X_n-a_n)$, for $a\in K^n$.
\end{theorem}

\begin{proof}
For $a\in K^n$, the evaluation $f\mapsto f(a)$ is a surjection onto the field $K$ whose kernel $\mathfrak m_a$ contains $X_i-a_i$; since every $f$ is $f(a)$ plus a combination of the $X_i-a_i$ (Taylor expansion at $a$),
$\mathfrak m_a=(X_1-a_1,\dots,X_n-a_n)$, and it is maximal by Volume~2, Proposition~6.8.

Let $\mathfrak m$ be a maximal ideal. Then $L=K[X_1,\dots,X_n]/\mathfrak m$ is a field, generated as a $K$-algebra by the classes $\bar X_i$, hence a finite extension of $K$ by Theorem~\ref{thm:zariskilemma}. As $K$ is algebraically closed,
$L=K$: each $\bar X_i$ equals some $a_i\in K$, that is $X_i-a_i\in\mathfrak m$. So $\mathfrak m_a\subseteq\mathfrak m$, and $\mathfrak m=\mathfrak m_a$ by maximality.

Finally, a proper ideal $I$ is contained in a maximal ideal $\mathfrak m_a$ (Theorem~\ref{thm:hilbertbasis} and the remark after Proposition~\ref{prop:noetherian}), and every $f\in I$ satisfies $f(a)=0$: $a\in V(I)$.
\end{proof}

The \emph{radical} of an ideal $J$ is $\sqrt J=\{f:f^m\in J\text{ for some }m\ge1\}$, an ideal containing $J$ (if $f^m,g^p\in J$, every term of $(f+g)^{m+p-1}$ is a multiple of $f^m$ or of $g^p$). A function that vanishes where $J$ vanishes need
not lie in $J$: $X$ vanishes on $V(X^2)=\{0\}$ but $X\notin(X^2)$. Powers are the only obstruction.

\begin{theorem}[Hilbert's Nullstellensatz, 1893]\label{thm:nullstellensatz}
Let $K$ be algebraically closed and $J$ an ideal of $K[X_1,\dots,X_n]$. Then $I(V(J))=\sqrt J$.
\end{theorem}

\begin{proof}
If $f^m\in J$ then $f^m$, hence $f$, vanishes on $V(J)$. Conversely let $f\in I(V(J))$, $f\ne0$, and $J=(g_1,\dots,g_k)$ (Theorem~\ref{thm:hilbertbasis}). The trick, due to J.~L.~Rabinowitsch in 1929, is to add a variable $Y$. The polynomials
$g_1,\dots,g_k,1-Yf$ have no common zero in $K^{n+1}$: where all $g_i$ vanish, so does $f$, and $1-Yf=1$. By Theorem~\ref{thm:weaknull} they generate the unit ideal:
\[1=\sum_ia_i(X,Y)g_i(X)+b(X,Y)\bigl(1-Yf(X)\bigr).\]
This is an identity of polynomials, so it remains true in the field of rational functions $K(X_1,\dots,X_n)$ after substituting $Y=1/f$, which kills the last term. Multiplying by a power $f^m$ with $m$ at least the degree of every
$a_i$ in $Y$ clears the denominators: $f^m=\sum_ic_i(X)g_i(X)$ with $c_i\in K[X_1,\dots,X_n]$. So $f\in\sqrt J$.
\end{proof}

\begin{corollary}[the dictionary]\label{cor:dictionary}
Over an algebraically closed field $K$, the maps $V$ and $I$ are inverse, inclusion-reversing bijections between radical ideals of $K[X_1,\dots,X_n]$ and closed subsets of $K^n$. Under them, prime ideals correspond to
irreducible closed sets, and maximal ideals to points.
\end{corollary}

\begin{proof}
$V(I(Y))=Y$ for closed $Y$ (Proposition~\ref{prop:zariski}), and $I(V(J))=\sqrt J=J$ for radical $J$ (Theorem~\ref{thm:nullstellensatz}); $I(Y)$ is radical since $f^m$ vanishes where $f$ does. Primes correspond to irreducible sets by
Proposition~\ref{prop:components}(b), and maximal ideals to points by Theorem~\ref{thm:weaknull}, since $V(\mathfrak m_a)=\{a\}$.
\end{proof}

\begin{think}
In $\C[X,Y]$, which ideal corresponds to the union of the two axes? To the origin counted ``twice'', $V(X,Y^2)$? Why does geometry not see the difference between $(X,Y)$ and $(X,Y^2)$, and why might one want it to?
\end{think}

The last question is the beginning of scheme theory, which keeps all ideals instead of only the radical ones. Chapter~4 takes the first step, by replacing the points of $K^n$ with the prime ideals of an arbitrary ring.

\section{Gröbner bases}

Hilbert's theorems say that ideals have finite bases and that a system without solutions has a certificate $1=\sum a_ig_i$. To compute, one needs a normal form: a way to decide whether $f\in(g_1,\dots,g_k)$. In one variable,
division with remainder by the gcd of the $g_i$ does it. In several variables, Buchberger's Gröbner bases do. We write $X^\alpha=X_1^{\alpha_1}\cdots X_n^{\alpha_n}$ for $\alpha\in\N^n$.

\begin{definition}\label{def:monomialorder}
A \emph{monomial order} is a total order $<$ on the monomials $X^\alpha$ such that $1\le X^\alpha$ for all $\alpha$, and $X^\alpha<X^\beta$ implies $X^{\alpha+\gamma}<X^{\beta+\gamma}$. For a non-zero polynomial $f$, its \emph{leading monomial}
$\operatorname{LM}(f)$ is the largest monomial occurring in it, $\operatorname{LC}(f)$ its coefficient and $\operatorname{LT}(f)=\operatorname{LC}(f)\operatorname{LM}(f)$ its leading term.
\end{definition}

The \emph{lexicographic order} compares exponents $\alpha,\beta$ at the first index where they differ ($X_1>X_2^{100}$); the \emph{graded reverse lexicographic order} compares total degrees first. An ideal generated by
monomials is a \emph{monomial ideal}; a polynomial lies in it if and only if each of its monomials is divisible by one of the generators, since a combination $\sum h_iX^{\alpha_i}$ only contains such monomials.

\begin{lemma}[Dickson]\label{lem:dickson}
A monomial ideal generated by a set $M$ of monomials is generated by a finite subset of $M$. Consequently a monomial order is a well-order: every non-empty set of monomials has a least element.
\end{lemma}

\begin{proof}
The ideal $(M)$ is finitely generated (Theorem~\ref{thm:hilbertbasis}), each generator being a combination of finitely many elements of $M$; these finitely many monomials generate. Given a non-empty set $M$ of monomials, let
$X^{\alpha_1},\dots,X^{\alpha_k}\in M$ generate $(M)$, and let $X^{\alpha_i}$ be the least of them. Every $X^\beta\in M$ is divisible by some $X^{\alpha_j}$, $X^\beta=X^{\alpha_j}X^\gamma\ge X^{\alpha_j}\ge X^{\alpha_i}$, using $1\le X^\gamma$ and
multiplicativity.
\end{proof}

\begin{proposition}[division]\label{prop:division}
Fix a monomial order and non-zero $g_1,\dots,g_k$. Every $f$ can be written $f=\sum_iq_ig_i+r$, where no monomial of $r$ is divisible by any $\operatorname{LM}(g_i)$, and where $\operatorname{LM}(q_ig_i)\le\operatorname{LM}(f)$ whenever $q_ig_i\ne0$.
\end{proposition}

\begin{proof}
Start with $p=f$, all $q_i=0$, $r=0$, and repeat while $p\ne0$: if $\operatorname{LM}(p)$ is divisible by some $\operatorname{LM}(g_i)$, take the first such $i$ and replace $p$ by $p-\frac{\operatorname{LT}(p)}{\operatorname{LT}(g_i)}g_i$, adding the
quotient to $q_i$; otherwise move $\operatorname{LT}(p)$ from $p$ to $r$. In both cases $f=\sum q_ig_i+r+p$ remains true, $\operatorname{LM}(p)$ strictly decreases, and the terms added to $q_ig_i$ have leading monomial $\operatorname{LM}(p)\le\operatorname{LM}(f)$.
A strictly decreasing sequence of monomials is finite by Lemma~\ref{lem:dickson}, so the process stops, with $p=0$.
\end{proof}

The remainder depends on the order of the $g_i$ in general, and $r\ne0$ does not prove $f\notin(g_1,\dots,g_k)$. Gröbner bases repair this.

\begin{definition}\label{def:groebner}
Let $I$ be a non-zero ideal. A finite set $G=\{g_1,\dots,g_k\}\subseteq I$ is a \emph{Gröbner basis} of $I$ if the leading monomials of $G$ generate the monomial ideal $(\operatorname{LM}(I))$ generated by all leading monomials of
elements of $I$.
\end{definition}

\begin{proposition}\label{prop:groebnerprops}
Every non-zero ideal has a Gröbner basis. If $G$ is a Gröbner basis of $I$, then $G$ generates $I$, the remainder of $f$ on division by $G$ does not depend on the choices made, and $f\in I$ if and only if this remainder is $0$.
\end{proposition}

\begin{proof}
By Lemma~\ref{lem:dickson}, $(\operatorname{LM}(I))$ is generated by finitely many $\operatorname{LM}(g_i)$, $g_i\in I$. Let $f\in I$ and divide: $f=\sum q_ig_i+r$ with $r\in I$. If $r\ne0$, then $\operatorname{LM}(r)\in(\operatorname{LM}(I))$ is divisible by some
$\operatorname{LM}(g_i)$, which the division forbids. So $r=0$ and $f\in(G)$. If $r,r'$ are two remainders of the same $f$, then $r-r'\in I$ and no monomial of $r-r'$ is divisible by any $\operatorname{LM}(g_i)$; as before $r-r'=0$. Finally
$f\in I$ if and only if the remainder, which lies in $f+I$, is $0$.
\end{proof}

To decide whether a given $G$ is a Gröbner basis, Buchberger looked at the cancellations between pairs. For non-zero $f,g$ with $X^\gamma=\operatorname{lcm}(\operatorname{LM}(f),\operatorname{LM}(g))$, the \emph{S-polynomial} is
\[S(f,g)=\frac{X^\gamma}{\operatorname{LT}(f)}f-\frac{X^\gamma}{\operatorname{LT}(g)}g,\]
built so that the leading terms cancel.

\begin{lemma}\label{lem:cancel}
Let $p=\sum_jc_jX^{\beta_j}g_{i_j}$, a finite sum with $c_j\in K$, in which every term $X^{\beta_j}g_{i_j}$ has the same leading monomial $X^\delta$, and suppose $\operatorname{LM}(p)<X^\delta$ (or $p=0$). Then $p$ is a linear
combination, with coefficients in $K$, of polynomials $X^{\delta-\gamma}S(g_{i_j},g_{i_l})$, where $X^\gamma$ is the lcm of the leading monomials of $g_{i_j}$ and $g_{i_l}$, each with leading monomial less than $X^\delta$.
\end{lemma}

\begin{proof}
Put $p_j=X^{\beta_j}g_{i_j}/\operatorname{LC}(g_{i_j})$, with leading term $X^\delta$, and $d_j=c_j\operatorname{LC}(g_{i_j})$, so $p=\sum_jd_jp_j$. The coefficient of $X^\delta$ in $p$ is $\sum_jd_j$, which must vanish. With $s$ terms, summation by parts gives
\[p=\sum_{j<s}(d_1+\dots+d_j)(p_j-p_{j+1})+(d_1+\dots+d_s)p_s=\sum_{j<s}(d_1+\dots+d_j)(p_j-p_{j+1}).\]
Each $p_j-p_{j+1}$ is $X^{\delta-\gamma}S(g_{i_j},g_{i_{j+1}})$, as one sees by factoring $X^{\delta-\gamma}$ (both $X^{\beta_j}\operatorname{LM}(g_{i_j})$ and $X^{\beta_{j+1}}\operatorname{LM}(g_{i_{j+1}})$ equal $X^\delta$, which is therefore a multiple of $X^\gamma$), or is $0$
when the two terms coincide; and its leading monomial is less than $X^\delta$, because the leading terms cancel.
\end{proof}

\begin{theorem}[Buchberger's criterion, 1965]\label{thm:buchberger}
A finite set $G=\{g_1,\dots,g_k\}$ of non-zero polynomials is a Gröbner basis of the ideal it generates if and only if, for all $i<j$, the remainder of $S(g_i,g_j)$ on division by $G$ is $0$.
\end{theorem}

\begin{proof}
If $G$ is a Gröbner basis, $S(g_i,g_j)\in(G)$ has remainder $0$ by Proposition~\ref{prop:groebnerprops}. Conversely, assume the condition, and let $f\in(G)$ be non-zero; we must show that $\operatorname{LM}(f)$ is divisible by some
$\operatorname{LM}(g_i)$. Among all expressions $f=\sum_ih_ig_i$, choose one for which $X^\delta=\max_i\operatorname{LM}(h_ig_i)$ is least (Lemma~\ref{lem:dickson}). Then $\operatorname{LM}(f)\le X^\delta$. If equality holds, $\operatorname{LM}(f)=\operatorname{LM}(h_i)\operatorname{LM}(g_i)$ for some
$i$, and we are done. Suppose $\operatorname{LM}(f)<X^\delta$, and split
\[f=\underbrace{\sum_{\operatorname{LM}(h_ig_i)=X^\delta}\operatorname{LT}(h_i)g_i}_{p}+\underbrace{\sum_{\operatorname{LM}(h_ig_i)=X^\delta}(h_i-\operatorname{LT}(h_i))g_i+\sum_{\operatorname{LM}(h_ig_i)<X^\delta}h_ig_i}_{\text{all terms }<X^\delta}.\]
So $\operatorname{LM}(p)<X^\delta$ or $p=0$, and Lemma~\ref{lem:cancel} writes $p$ as a combination of $X^{\delta-\gamma}S(g_i,g_j)$. By hypothesis and Proposition~\ref{prop:division}, $S(g_i,g_j)=\sum_la_{ijl}g_l$ with
$\operatorname{LM}(a_{ijl}g_l)\le\operatorname{LM}(S(g_i,g_j))$, so $X^{\delta-\gamma}S(g_i,g_j)=\sum_lX^{\delta-\gamma}a_{ijl}g_l$ with every term of leading monomial less than $X^\delta$. Substituting gives an expression of $f$ whose terms all have leading
monomial less than $X^\delta$, contradicting the choice of $X^\delta$.
\end{proof}

\begin{corollary}[Buchberger's algorithm]\label{cor:buchalg}
Start from generators $G$ of $I$. As long as some $S(g_i,g_j)$ has a non-zero remainder $r$ on division by $G$, add $r$ to $G$. This process stops, and then $G$ is a Gröbner basis of $I$.
\end{corollary}

\begin{proof}
Each added $r$ lies in $I$, and $\operatorname{LM}(r)$ is not divisible by any $\operatorname{LM}(g)$, $g\in G$, so the monomial ideal $(\operatorname{LM}(G))$ grows strictly. By the ascending chain condition this happens finitely often, and at
the end Theorem~\ref{thm:buchberger} applies.
\end{proof}

For the lexicographic order, a Gröbner basis is triangular, like Gaussian elimination: its last elements involve only the last variable, and the system can be solved one variable at a time (Exercise~\ref{ex:groebnerhand}).
The weak Nullstellensatz becomes an algorithm: over an algebraically closed field, a system has no solution if and only if its reduced Gröbner basis is $\{1\}$.

\begin{beauty}
Buchberger named his bases after his thesis adviser at Innsbruck, Wolfgang Gröbner. The same idea had appeared a year earlier, for power series and under the name of
\emph{standard bases}, in Heisuke Hironaka's work on the resolution of singularities. But it was Buchberger's algorithm that made it a tool, and today every computer algebra system contains it. It designs robot arms by solving their inverse
kinematics (Volume~8, Chapter~1), proves theorems of elementary geometry automatically by checking ideal membership, and decides whether a graph can be coloured with three colours or a Sudoku solved
(Exercise~\ref{ex:colouring}). Gordan's theology became computation.
\end{beauty}

\section{Exercises}

\begin{exercise}[$\star$]\label{ex:nonnoeth}
Show that the subring $K[X,XY,XY^2,XY^3,\dots]$ of the Noetherian ring $K[X,Y]$ is not Noetherian.
\end{exercise}

\begin{exercise}[$\star$]\label{ex:radicals}
Over $\C$, find the irreducible components of $V(X^2-Y^2)$ and of $V(XY,XZ,YZ)\subseteq\C^3$, and compute $I(V(X^2Y,XY^2))$ in $\C[X,Y]$.
\end{exercise}

\begin{exercise}[$\star\star$]\label{ex:notproduct}
Let $K$ be infinite. Show that the diagonal $\{(x,x)\}$ is closed in the Zariski topology of $K^2$, but not in the product of the Zariski topologies of $K$ and $K$. Deduce that the Zariski topology of $K^2$ is strictly finer
than the product topology.
\end{exercise}

\begin{exercise}[$\star\star$]\label{ex:realnull}
Over $\R$: show that every algebraic subset of $\R^n$ is the zero set of a single polynomial, and that the Nullstellensatz fails. Which ideal $J$ of $\R[X,Y]$ has $I(V(J))\ne\sqrt J$ with $V(J)$ a single point?
\end{exercise}

\begin{exercise}[$\star\star$]\label{ex:colouring}
Let $\Gamma$ be a finite graph with vertex set $\{1,\dots,n\}$, and $\omega=e^{2\pi i/3}$. Show that $\Gamma$ can be properly coloured with three colours if and only if the polynomials $X_v^3-1$ ($v$ a vertex) and
$X_u^2+X_uX_v+X_v^2$ ($uv$ an edge) have a common zero in $\C^n$, if and only if they do not generate the unit ideal of $\C[X_1,\dots,X_n]$.
\end{exercise}

\begin{exercise}[$\star\star$]\label{ex:groebnerhand}
Compute by hand, with Buchberger's algorithm, a Gröbner basis of $I=(X^2-Y,\,X^3-X)$ in $\C[X,Y]$ for the lexicographic order with $X>Y$, and use it to find $V(I)$. Is $I$ radical?
\end{exercise}

\section{Hints}

\begin{hint}{ex:nonnoeth}
Consider the chain of ideals generated by $XY,\dots,XY^k$.
\end{hint}

\begin{hint}{ex:radicals}
Factor. For the last one, the zero set is the union of the axes.
\end{hint}

\begin{hint}{ex:notproduct}
In the product of two cofinite topologies, a non-empty open set contains a product of two cofinite sets.
\end{hint}

\begin{hint}{ex:realnull}
Sums of squares. For the last question, think of $X^2+Y^2$.
\end{hint}

\begin{hint}{ex:colouring}
The roots of $X^3-1$ are $1,\omega,\omega^2$, and for two of them $x^3-y^3=(x-y)(x^2+xy+y^2)$.
\end{hint}

\begin{hint}{ex:groebnerhand}
Start with $S(X^3-X,X^2-Y)$; you should find three polynomials in the end.
\end{hint}

\section{Solutions}

\begin{solution}{ex:nonnoeth}
\idea The generators $XY^k$ cannot be expressed through each other inside the subring, because the subring has no $Y$ alone.
\writeup Let $S$ be the subring: it is spanned by $1$ and the monomials $X^aY^b$ with $a\ge1$. Let $J_k=(XY,\dots,XY^k)$, an ideal of $S$, whose elements are the sums $\sum_{j\le k}s_jXY^j$ with $s_j\in S$. A monomial of $s_jXY^j$ is
either $cXY^j$, from the constant term of $s_j$, or $X^{a+1}Y^{b+j}$ with $a\ge1$, of degree at least $2$ in $X$. So the only monomials of degree $1$ in $X$ occurring in elements of $J_k$ are the $XY^j$ with $j\le k$, and
$XY^{k+1}\notin J_k$. Hence $J_1\subsetneq J_2\subsetneq\cdots$, and $S$ is not Noetherian. Subrings of Noetherian rings need not be Noetherian; quotients are.
\end{solution}

\begin{solution}{ex:radicals}
\idea Factor into irreducibles, which give prime principal ideals in a unique factorisation domain.
\writeup $X^2-Y^2=(X-Y)(X+Y)$, so $V(X^2-Y^2)$ is the union of the lines $Y=X$ and $Y=-X$, both irreducible (their ideals $(X\mp Y)$ are prime, the quotient being $\C[X]$). A point of $V(XY,XZ,YZ)$ has at least two zero
coordinates, so the set is the union of the three coordinate axes, each irreducible and none containing another. Finally $V(X^2Y,XY^2)=V(XY)$ is the union of the axes, so $I(V(X^2Y,XY^2))=\sqrt{(X^2Y,XY^2)}=(XY)$: indeed $(XY)^2=X\cdot XY^2$
lies in the ideal, and a polynomial vanishing on both axes is divisible by $X$ and by $Y$.
\end{solution}

\begin{solution}{ex:notproduct}
\idea A non-empty product open set is huge, and cannot avoid the diagonal.
\writeup The diagonal is $V(X-Y)$. In the product topology, the closed sets of $K$ are the finite sets and $K$, and a basic open set is $U_1\times U_2$ with $U_i$ cofinite. If the diagonal $\Delta$ were closed, its complement would contain
some non-empty $U_1\times U_2$; but $U_1\cap U_2$ is cofinite, hence non-empty as $K$ is infinite, and a point $(x,x)$ with $x\in U_1\cap U_2$ lies in $U_1\times U_2$ and in $\Delta$. The product topology is coarser than the Zariski
topology (the projections are polynomial, hence continuous for the Zariski topology, since the preimage of $V(f)$ is $V(f(X))$), and the diagonal shows it is strictly coarser.
\end{solution}

\begin{solution}{ex:realnull}
\idea Over $\R$, $\sum f_i^2=0$ if and only if every $f_i=0$.
\writeup $V(f_1,\dots,f_k)=V(f_1^2+\dots+f_k^2)$ in $\R^n$, and every algebraic set is of the form $V(f_1,\dots,f_k)$ by Corollary~\ref{cor:finitesystem}. For $J=(X^2+Y^2)$, $V(J)=\{(0,0)\}$, so $I(V(J))=(X,Y)$, while $X\notin\sqrt J$:
the polynomials $X^m$ are not multiples of $X^2+Y^2$ (set $X=1$, $Y=i$: $1^m\ne0$ but $X^2+Y^2$ vanishes). And $J=(X^2+1)$ in $\R[X]$ is proper with $V(J)=\varnothing$.
\end{solution}

\begin{solution}{ex:colouring}
\idea Colours are the three cube roots of unity, and the edge polynomial is what $x^3-y^3$ becomes after removing the factor $x-y$.
\writeup The common zeros of the $X_v^3-1$ are the maps $x\colon V\to\{1,\omega,\omega^2\}$, colourings with three colours. For cube roots of unity $x,y$, $0=x^3-y^3=(x-y)(x^2+xy+y^2)$, so if $x\ne y$ then $x^2+xy+y^2=0$;
and if $x=y$ then $x^2+xy+y^2=3x^2\ne0$. So the edge polynomials vanish exactly when adjacent vertices get different colours: common zeros are proper colourings. By the weak Nullstellensatz (Theorem~\ref{thm:weaknull}) a common zero
exists if and only if the ideal is proper, that is, if and only if $1$ is not in it, and a Gröbner basis computation decides this.
\end{solution}

\begin{solution}{ex:groebnerhand}
\idea Run the algorithm; the lexicographic order eliminates $X$ as far as it can.
\writeup Let $g_1=X^2-Y$, $g_2=X^3-X$. $S(g_2,g_1)=g_2-Xg_1=XY-X$; its leading monomial $XY$ is not divisible by $X^2$ or $X^3$, so add $g_3=XY-X$. Now $S(g_1,g_3)=Yg_1-Xg_3=-Y^2+X^2$, which reduces by $g_1$ to $Y-Y^2$; add
$g_4=Y^2-Y$. The remaining S-polynomials reduce to $0$: $S(g_3,g_4)=Yg_3-Xg_4=-XY+XY=0$, $S(g_1,g_4)=Y^2g_1-X^2g_4=X^2Y-Y^3$, which reduces by $Yg_1$ to $Y^2-Y^3=-Yg_4$, then to $0$; and $g_2=Xg_1+g_3$ is now redundant,
since $\operatorname{LM}(g_2)=X^3$ is a multiple of $X^2$. So $\{X^2-Y,\,XY-X,\,Y^2-Y\}$ is a Gröbner basis. From $Y^2=Y$, $Y\in\{0,1\}$. If $Y=0$: $X^2=0$ and $X\cdot0=X$ gives $X=0$. If $Y=1$: $X^2=1$, $X=\pm1$. So $V(I)=\{(0,0),(1,1),(-1,1)\}$. The
ideal is radical: the quotient $\C[X,Y]/I$ has the basis $1,X,Y$ given by the monomials not divisible by any leading monomial, so it has dimension $3$, as does the ring of functions on three points, onto which it maps
surjectively (interpolation); hence it is isomorphic to it, with no nilpotent elements.
\end{solution}

\section*{Sources and further reading}
\addcontentsline{toc}{section}{Sources and further reading}

D.~Hilbert, ``Über die Theorie der algebraischen Formen'', \emph{Math.\ Ann.} 36 (1890), 473--534, and ``Über die vollen Invariantensysteme'', \emph{Math.\ Ann.} 42 (1893), 313--373; P.~Gordan, ``Beweis, dass jede Covariante und
Invariante einer binären Form eine ganze Function mit numerischen Coefficienten einer endlichen Anzahl solcher Formen ist'', \emph{J.\ Reine Angew.\ Math.} 69 (1868), 323--354. Gordan's remark and its sequel are told in
C.~Reid, \emph{Hilbert} (Springer, 1970), Chapter~5, and retold in ``We must know, we will know'', \emph{Plus Magazine} (Cambridge). E.~Noether, ``Idealtheorie in Ringbereichen'', \emph{Math.\ Ann.} 83 (1921), 24--66;
J.~L.~Rabinowitsch, ``Zum Hilbertschen Nullstellensatz'', \emph{Math.\ Ann.} 102 (1929), 520; O.~Zariski, ``A new proof of Hilbert's Nullstellensatz'', \emph{Bull.\ Amer.\ Math.\ Soc.} 53 (1947), 362--368. The proof of
Zariski's lemma through the Artin--Tate lemma follows Atiyah and Macdonald, Chapter~7. B.~Buchberger, \emph{Ein Algorithmus zum Auffinden der Basiselemente des Restklassenringes nach einem nulldimensionalen Polynomideal},
thesis, Innsbruck, 1965 (English translation in \emph{J.\ Symbolic Comput.} 41 (2006), 475--511).

For more: D.~Cox, J.~Little and D.~O'Shea, \emph{Ideals, Varieties, and Algorithms} (Springer, 4th ed.\ 2015), the best introduction to the subject, with Chapters~1--4 covering this chapter in detail; D.~Eisenbud,
\emph{Commutative Algebra with a View Toward Algebraic Geometry} (Springer, 1995), Chapters~1, 4 and 15.
